\documentclass[twoside, 10pt]{article}
\usepackage[utf8]{inputenc}
\usepackage{amsmath}
\usepackage{amssymb}
\usepackage{amsfonts}
\usepackage{geometry}
\usepackage{booktabs}
\usepackage{graphicx}
\usepackage{float}
\usepackage{xcolor}
\usepackage{enumitem}
\usepackage{caption}
\usepackage{tabularx}
\geometry{margin=1in}
\usepackage{hyperref}
\usepackage{listings}

% Page geometry
% (duplicate geometry statement removed)

% Color definitions
\definecolor{primaryblue}{RGB}{26, 54, 93}
\definecolor{darkgrey}{RGB}{45, 55, 72}

% Global paragraph settings
\setlength{\parindent}{0pt}
\setlength{\parskip}{0.8em}
\raggedbottom

% --- FIXED NUMBERING MACROS ---
\renewcommand{\thesubsection}{\arabic{subsection}}
\renewcommand{\thesubsubsection}{\thesubsection.\arabic{subsubsection}} 

% Define the paragraph macro to act as your 3-digit sub-subsubsection (2.4.1)
\renewcommand{\theparagraph}{\thesubsubsection.\arabic{paragraph}} 

% Force paragraph headings to have their own standalone line instead of running inline
\makeatletter
\renewcommand{\paragraph}{\@startsection{paragraph}{4}{0pt}%
    {-3.25ex plus -1ex minus -0.2ex}%
    {1.5ex plus 0.2ex}%
    {\normalfont\normalsize\bfseries}}
\makeatother

\usepackage{titlesec}

% Ensure LaTeX numbers headings down to the paragraph level
\setcounter{secnumdepth}{4}
% ------------------------------

% Modify subsection: keep default size (\large) and alignment, but force normal font weight
\titleformat{\subsection}
  {\normalfont\large\fontseries{m}\selectfont} % Medium weight instead of bold
  {\thesubsection.}                            % Format for the number (e.g., 1.)
  {0.5em}                                      % Spacing between number and text
  {}

  % The format is: \titlespacing*{command}{left spacing}{space before}{space after}
\titlespacing*{\section}{0pt}{2ex plus 1ex minus .2ex}{0.5ex plus .2ex}
\titlespacing*{\subsection}{0pt}{1ex plus 1ex minus .2ex}{0.5ex plus .2ex}
\titlespacing*{\subsubsection}{0pt}{0.5ex plus 0.5ex minus .2ex}{0.2ex plus .2ex}

% Setup for Python code formatting
\lstset{
    language=Python,
    basicstyle=\ttfamily\small,
    keywordstyle=\color{blue},
    stringstyle=\color{red},
    commentstyle=\color{green!60!black},
    numbers=left,
    numberstyle=\tiny\color{gray},
    stepnumber=1,
    breaklines=true,
    frame=single,
    showstringspaces=false
}

\begin{document}

\title{\textbf{\Large HomeWork4}}
\author{Aswath Oruganti}
\maketitle

\section{Comparing Classifiers: Divorce Classification/Prediction}
\subsection{Classifier Performance}
The classifiers were trained on the 54-dimensional feature space. For the Gaussian Naive Bayes classifier, the variance smoothing parameter was set to $\epsilon = 10^{-3}$ to prevent errors associated with zero-variance features. The testing accuracies are reported in Table \ref{tab:accuracies}.
\begin{table}[h]
    \centering
    \caption{Testing Accuracies of Classifiers (54 Features)}\label{tab:accuracies}
    \vspace{0.1in}
    \begin{tabular}{lc}
        \toprule
        \textbf{Classifier} & \textbf{Testing Accuracy} \\
        \midrule
        Naive Bayes (Gaussian) & 97.06\% \\
        Logistic Regression & 97.06\% \\
        K-Nearest Neighbors ($k=5$) & 97.06\% \\
        \bottomrule
    \end{tabular}
\end{table}
\subsection*{Performance Comparison:}
All three classifiers perform remarkably well—and identically—on the test set, achieving an accuracy of 97.06\%.  
The underlying reason for this exceptional performance is the nature of the dataset. we know the data is polarized by looking at the raw CSV file and understanding the psychological nature of the survey. 
The UCI dataset consists of 54 statements (e.g., ``We don't have time to communicate as partners'', ``I know my spouse's basic anxieties''). Participants rate these on a 5-point scale (0 to 4).
\begin{description}
    \item[Extreme Correlation:] Human behavior in failing marriages tends to cascade. If a couple scores a `4' (strongly agree) on poor communication, they almost always score `4's on lack of trust, frequent arguing, and feeling ignored.
    \item[The Polarization:] Consequently, the dataset does not have many people scoring a ``mixed bag'' of 2s and 3s. Happily married couples answer mostly 0s and 1s across the board, while divorcing couples answer mostly 3s and 4s across the board.
\end{description}
Because the answers move together in bulk, the data points naturally herd into two tight, opposing corners of the feature space, creating those distinct, easily separable clusters we observed.
Because of this polarization, the dataset is nearly \textbf{linearly separable} with two distinct, tight clusters in the feature space. 
\begin{itemize}
    \item \textbf{Logistic Regression} easily finds a linear hyperplane separating these distant clusters.
    \item \textbf{KNN} performs perfectly because data points of the same class are tightly packed together.
    \item \textbf{Naive Bayes} succeeds because the strong conditional probabilities of the features easily outweigh the violated "conditional independence" assumption.
\end{itemize}

\subsection{PCA and Decision Boundaries}
To visualize the decision boundaries, Principal Component Analysis (PCA) was performed to project the 54-dimensional data into a 2-dimensional space. The three classifiers were subsequently retrained on these two principal components. The resulting decision boundaries and actual data points are illustrated in Figure \ref{fig:boundaries}.
\begin{figure}[h]
    \centering
    % Ensure decision_boundaries.png is in the same directory as this .tex file
    \includegraphics[width=\textwidth]{ISYE6740_Summer_2026_HW4/decision_boundaries.png}
    \caption{Decision boundaries of Naive Bayes, Logistic Regression, and KNN on the 2D PCA-reduced dataset. Red regions indicate Class 0 predictions, while blue regions indicate Class 1 predictions.}\label{fig:boundaries}
\end{figure}
\subsection*{Decision Boundaries Discussion:}
By projecting the data into 2D space, the fundamental geometric differences between how these algorithms classify data become visually apparent:
\begin{itemize}
    \item \textbf{Naive Bayes (Quadratic Boundary):} The boundary is distinctly curved. Because Gaussian Naive Bayes fits a multivariate normal distribution (mean and variance) to each class separately, the intersection of these probability distributions forms a quadratic surface (a curve in 2D).
    \item \textbf{Logistic Regression (Linear Boundary):} The boundary is strictly a straight line. Logistic regression inherently models the log-odds as a linear combination of the input features, forcing the decision boundary to always be linear regardless of the data's true underlying shape.
    \item \textbf{K-Nearest Neighbors (Non-linear/Local Boundary):} The boundary is highly irregular and jagged. KNN does not assume any global mathematical distribution or equation. Instead, it adapts locally based on the proximity of the $k$ nearest training points, resulting in isolated ``islands'' and complex edges that perfectly wrap around the local data clusters.
\end{itemize}

\section{Comparing mult-class classifiers for MNIST Dataset}
\subsection{Classifier Performance Metrics}
Based on the classification results, bar charts, and knee plots provided, the following section synthesizes the performance of the five classifiers (KNN, Logistic Regression, Linear SVM, Kernel SVM, and MLP) across both the MNIST Digits and MNIST Fashion datasets.
\subsubsection{Tuning Methodology (KNN)}
To find the optimal $K$ for the \texttt{KNeighborsClassifier}, validation accuracy was evaluated against the number of neighbors.
\begin{itemize}
    \item \textbf{MNIST Digits:} The knee plot demonstrates that validation accuracy is strictly decreasing as $K$ increases. The highest accuracy is achieved at \textbf{K = 1}.
    \item \textbf{MNIST Fashion:} The knee plot shows a peak in validation accuracy at \textbf{K = 5}, after which performance steadily declines.
\end{itemize}
\begin{figure}[H]
    \centering
    \begin{minipage}{0.48\textwidth}
        \centering
        \includegraphics[width=\linewidth]{{"ISYE6740_Summer_2026_HW4/knn_knee_plot_MNIST Digits.png"}}
        \caption{KNN Knee Plot for MNIST Digits}
    \end{minipage}\hfill
    \begin{minipage}{0.48\textwidth}
        \centering
        \includegraphics[width=\linewidth]{{"ISYE6740_Summer_2026_HW4/knn_knee_plot_MNIST Fashion.png"}}
        \caption{KNN Knee Plot for MNIST Fashion}
    \end{minipage}
\end{figure}
\subsubsection{Metric Summaries: MNIST Digits}
This dataset proved highly solvable across all models, with most classifiers achieving $>90\%$ on Precision, Recall, and F1-scores across all 10 digit classes.
\begin{table}[H]
    \centering
    \begin{tabularx}{\textwidth}{l l X l}
        \toprule
        \textbf{Classifier} & \textbf{Best Parameter} & \textbf{General Performance Trend} & \textbf{Weakest Class} \\
        \midrule
        Kernel SVM (RBF) & C = 1.0 ($m=5000$) & Excellent (F1 $> 0.95$ across all classes). & Class 8 \\
        KNN & K = 1 & Excellent (F1 $> 0.95$ across all classes). & Class 5 \& 8 \\
        MLP & (20, 10) & Very Good (F1 $\approx 0.92$ to $0.98$). & Class 3 \& 8 \\
        Linear SVM & C = 0.1 & Good (F1 $> 0.88$ across all classes). & Class 5 \& 8 \\
        Logistic Regression & Default & Good (F1 $> 0.88$ across all classes). & Class 5 \& 8 \\
        \bottomrule
    \end{tabularx}
    \caption{Performance Summary for MNIST Digits}
\end{table}
\noindent Across most models, classes 5 and 8 consistently showed the lowest F1-scores, indicating slight confusion in linearly separating these specific digits.
\begin{figure}[H]
    \centering
    \begin{minipage}{0.48\textwidth}
        \centering
        \includegraphics[width=\linewidth]{{"ISYE6740_Summer_2026_HW4/Kernel SVM RBF (C=1.0)_plot_MNIST Digits.png"}}
        \caption{Kernel SVM on Digits}
    \end{minipage}\hfill
    \begin{minipage}{0.48\textwidth}
        \centering
        \includegraphics[width=\linewidth]{{"ISYE6740_Summer_2026_HW4/MLP (20, 10)_plot_MNIST Digits.png"}}
        \caption{MLP on Digits}
    \end{minipage}
    \vspace{1em}
    
    \begin{minipage}{0.48\textwidth}
        \centering
        \includegraphics[width=\linewidth]{{"ISYE6740_Summer_2026_HW4/KNN (K=1)_plot_MNIST Digits.png"}}
        \caption{KNN on Digits}
    \end{minipage}\hfill
    \begin{minipage}{0.48\textwidth}
        \centering
        \includegraphics[width=\linewidth]{{"ISYE6740_Summer_2026_HW4/Logistic Regression_plot_MNIST Digits.png"}}
        \caption{Log. Regression on Digits}
    \end{minipage}
\end{figure}
\subsubsection{Metric Summaries: MNIST Fashion}
This dataset was significantly harder for all classifiers. The inter-class variance is much higher, leading to wider disparities in class-by-class performance.
\begin{table}[H]
    \centering
    \begin{tabularx}{\textwidth}{l l X l}
        \toprule
        \textbf{Classifier} & \textbf{Best Parameter} & \textbf{General Performance Trend} & \textbf{Weakest Class} \\
        \midrule
        Kernel SVM (RBF) & C = 10.0 ($m=5000$) & Best overall. F1 ranges from $\approx 0.68$ to $0.98$. & Class 6 (F1 $\approx 0.68$) \\
        MLP & (20, 10) & Strong contender. F1 ranges from $\approx 0.65$ to $0.97$. & Class 6 (F1 $\approx 0.65$) \\
        KNN & K = 5 & Moderate to Good. F1 ranges from $\approx 0.60$ to $0.95$. & Class 6 (F1 $\approx 0.62$) \\
        Linear SVM & C = 0.1 & Moderate. F1 ranges from $\approx 0.60$ to $0.95$. & Class 6 (F1 $\approx 0.62$) \\
        Logistic Regression & Default & Moderate. F1 ranges from $\approx 0.60$ to $0.95$. & Class 6 (F1 $\approx 0.62$) \\
        \bottomrule
    \end{tabularx}
    \caption{Performance Summary for MNIST Fashion}
\end{table}
\noindent Class 6 (``Shirt'') is universally the lowest-performing class across every single classifier, including the MLP. It suffers from terrible recall, meaning actual shirts are frequently misclassified as other tops (T-shirts, pullovers, or coats). Class 1 (Trouser) is universally the easiest to classify.
\begin{figure}[H]
    \centering
    \begin{minipage}{0.48\textwidth}
        \centering
        \includegraphics[width=\linewidth]{{"ISYE6740_Summer_2026_HW4/Kernel SVM RBF (C=10.0)_plot_MNIST Fashion.png"}}
        \caption{Kernel SVM on Fashion}
    \end{minipage}\hfill
    \begin{minipage}{0.48\textwidth}
        \centering
        \includegraphics[width=\linewidth]{{"ISYE6740_Summer_2026_HW4/MLP (20, 10)_plot_MNIST Fashion.png"}}
        \caption{MLP on Fashion}
    \end{minipage}
    \vspace{1em}
    
    \begin{minipage}{0.48\textwidth}
        \centering
        \includegraphics[width=\linewidth]{{"ISYE6740_Summer_2026_HW4/KNN (K=5)_plot_MNIST Fashion.png"}}
        \caption{KNN on Fashion}
    \end{minipage}\hfill
    \begin{minipage}{0.48\textwidth}
        \centering
        \includegraphics[width=\linewidth]{{"ISYE6740_Summer_2026_HW4/Linear SVM (C=0.1)_plot_MNIST Fashion.png"}}
        \caption{Linear SVM on Fashion}
    \end{minipage}
\end{figure}
\subsection{Discussion of Performance}
\subsubsection{Comparing Classifiers Against Each Other}
When comparing the five models, a clear hierarchy of capability emerges based on their ability to handle non-linear decision boundaries:
\begin{enumerate}
    \item \textbf{Kernel SVM (RBF):} This is the top performer. The RBF Kernel projects the data into a higher-dimensional space where complex overlapping classes can be separated. Even with only a 5,000-sample subset, it outperforms the linear models trained on the full dataset, particularly on the harder Fashion dataset.
    \item \textbf{MLP (20, 10) \& KNN:} These follow closely behind.
    \begin{itemize}
        \item The \textbf{MLP} serves as a strong non-linear baseline. Even with a relatively small architecture, it learns hierarchical features that allow it to beat linear models comfortably on the Fashion dataset.
        \item \textbf{KNN} performs exceptionally well on Digits ($K=1$) because the Euclidean distance between centralized, stark black-and-white images is a surprisingly strong baseline for similarity.
    \end{itemize}
    \item \textbf{Linear SVM \& Logistic Regression:} These models performed the worst overall. They are strictly limited to linear decision boundaries. In high-dimensional pixel space (784 dimensions), drawing a single flat hyperplane to separate visually similar classes results in severe overlap and lower accuracy.
\end{enumerate}
\subsubsection{Insight into Different Datasets (Digits vs. Fashion)}
There is a massive performance gap between the two datasets, illustrating why machine learning benchmarks largely transitioned from Digits to Fashion.
\begin{itemize}
    \item \textbf{MNIST Digits} is essentially a solved problem for basic machine learning. The images are stark black and white strokes. Because the intra-class variance is small compared to the inter-class variance, even rigid linear models perform wonderfully.
    \item \textbf{MNIST Fashion} is inherently more complex. It consists of grayscale photos of clothing with high intra-class variance (e.g., a ``dress'' with different sleeve lengths) and extremely high inter-class similarity. Classes like Shirts, T-shirts, Pullovers, and Coats all share the same basic geometric silhouette.
\end{itemize}
\textbf{Conclusion:} Models capable of capturing non-linear relationships (Kernel SVM and MLP) are strictly necessary for datasets with high inter-class similarity like MNIST Fashion. Linear models fail to capture the nuanced pixel interactions required to differentiate a long-sleeve shirt from a light jacket.
\subsubsection{Appendix: Understanding the F1-Score}
The \textbf{F1-score} is a performance metric that combines both Precision and Recall into a single, balanced number using the harmonic mean:
\begin{equation}
F1 = 2 \times \frac{\text{Precision} \times \text{Recall}}{\text{Precision} + \text{Recall}}
\end{equation}
The harmonic mean heavily penalizes extreme imbalances. If a model has high Precision but terrible Recall (e.g., it only guesses ``Shirt'' when absolutely certain, but misses 80\% of actual shirts), the F1-score will drop significantly. To achieve a high F1-score, a model must successfully balance both quality (Precision) and quantity (Recall).

\section{Support Vector Machine (SVM)}
\subsection{\texorpdfstring{Why We Can Set the Margin c=1 in the SVM Formulation}{Why We Can Set the Margin c=1 in the SVM Formulation}}
\subsubsection{Derivation: Distance Between the Two Hyperplanes}
To mathematically prove why the distance between the two margin hyperplanes is $\frac{2c}{||w||}$, let's denote the two parallel hyperplanes as $H_+$ and $H_-$:
\begin{align*}
    H_+ &: w^T x + b = c \\
    H_- &: w^T x + b = -c
\end{align*}
The weight vector $w$ is orthogonal (normal) to both hyperplanes. Therefore, the unit normal vector pointing directly from one hyperplane to the other is:
\[ \hat{w} = \frac{w}{||w||} \]
Let $x_+$ be any arbitrary point lying on the positive hyperplane $H_+$, and $x_-$ be any arbitrary point lying on the negative hyperplane $H_-$. By definition, these points must satisfy their respective hyperplane equations:
\begin{align}
    w^T x_+ + b &= c \implies w^T x_+ = c - b \\
    w^T x_- + b &= -c \implies w^T x_- = -c - b
\end{align}
The shortest geometric distance $D$ between two parallel planes is the length of the projection of any vector connecting them (like $x_+ - x_-$) onto the normal vector $\hat{w}$:
\[ D = \left| (x_+ - x_-) \cdot \hat{w} \right| = \left| (x_+ - x_-)^T \frac{w}{||w||} \right| \]
Using the distributive property of the dot product, we can expand the numerator:
\[ D = \frac{|w^T x_+ - w^T x_-|}{||w||} \]
Now, substitute the expressions from equations (1) and (2) into the numerator:
\[ D = \frac{|(c - b) - (-c - b)|}{||w||} \]
\[ D = \frac{|c - b + c + b|}{||w||} \]
\[ D = \frac{|2c|}{||w||} \]
Since the margin parameter $c$ is strictly positive, the absolute value is redundant. 
\[ D = \frac{2c}{||w||} \]
Thus, the geometric distance between the two supporting hyperplanes is exactly $\frac{2c}{||w||}$. 
\subsubsection{Proof and Justification:} 
The SVM decision boundary is a hyperplane defined by the equation $w^T x + b = 0$. The geometric margin (the distance from the hyperplane to the closest data points) is determined by the parallel hyperplanes supporting the two classes:
\begin{align*}
    w^T x + b &= c \quad \text{(for positive class)} \\
    w^T x + b &= -c \quad \text{(for negative class)}
\end{align*}
The distance between these two supporting hyperplanes is $\frac{2c}{||w||}$ as shown in subsubsection 3.1.1.
Because $w$ and $b$ define a hyperplane, their scale is arbitrary. If we multiply the weights and bias by any non-zero positive scalar $\kappa$, the decision boundary $w^T x + b = 0$ remains exactly the same, but our margin equations become:
\[ \kappa w^T x + \kappa b = \kappa c \]
To simplify our optimization problem, we can choose the scalar $\kappa = \frac{1}{c}$. By doing this, we re-scale our parameters such that:
\[ w_{new} = \frac{w}{c}, \quad b_{new} = \frac{b}{c} \]
This forces the margin boundaries to be precisely $1$ and $-1$. Setting $c=1$ simply removes this scale invariance, fixing the functional margin to $1$ and reducing the problem to maximizing $\frac{2}{||w||}$ (or minimizing $\frac{1}{2}||w||^2$) subject to the constraint $y_i(w^T x_i + b) \ge 1$.

\subsection{Lagrangian Dual Formulation}
\subsubsection{The Primal Optimization Problem}
The primal optimization problem for a hard-margin SVM seeks to maximize the geometric margin by minimizing the norm of the weight vector, subject to the constraint that all data points are correctly classified with a functional margin of at least 1.
\begin{itemize}
    \item \textbf{Minimize:} $\frac{1}{2}||w||^2$
    \item \textbf{Subject to:} $y_i(w^T x_i + b) \ge 1 \quad \forall i \in \{1, \dots, n\}$
\end{itemize}
\subsubsection{The Lagrangian Formulation}
We rewrite the constraints as $1 - y_i(w^T x_i + b) \le 0$ and introduce the Lagrange multipliers (dual variables) $\alpha_i \ge 0$. The Lagrangian integrates the objective function and the constraints:
\begin{equation}
    \mathcal{L}(w, b, \alpha) = \frac{1}{2}||w||^2 + \sum_{i=1}^n \alpha_i \left( 1 - y_i(w^T x_i + b) \right)
\end{equation}
\subsubsection{Deriving the Dual Formulation}
To find the dual, we minimize $\mathcal{L}$ with respect to the primal variables $w$ and $b$. We take the partial derivatives and set them to zero.
\textbf{Derivative with respect to $w$:}
\begin{equation}
    \frac{\partial \mathcal{L}}{\partial w} = w - \sum_{i=1}^n \alpha_i y_i x_i = 0 \implies w = \sum_{i=1}^n \alpha_i y_i x_i
\end{equation}
\textbf{Derivative with respect to $b$:}
\begin{equation}
    \frac{\partial \mathcal{L}}{\partial b} = -\sum_{i=1}^n \alpha_i y_i = 0 \implies \sum_{i=1}^n \alpha_i y_i = 0
\end{equation}
Substituting equations (2) and (3) back into the original Lagrangian yields the Dual formulation, which depends strictly on the dual variables $\alpha$ and the dot products of the data points:
\begin{equation}
    \mathcal{L}_D(\alpha) = \sum_{i=1}^n \alpha_i - \frac{1}{2} \sum_{i=1}^n \sum_{j=1}^n \alpha_i \alpha_j y_i y_j (x_i^T x_j)
\end{equation}
\subsubsection{Crucial Implications}
The derivation $w = \sum_{i=1}^n \alpha_i y_i x_i$ is profoundly important in machine learning theory:
\begin{enumerate}
    \item \textbf{Sparsity and Support Vectors:} Due to KKT complementary slackness, $\alpha_i > 0$ only for data points that lie exactly on the margin. For all other points, $\alpha_i = 0$. Therefore, the optimal weight vector $w$ is merely a linear combination of a small subset of the training data points (the Support Vectors).
    \item \textbf{Inner Product Dependency:} The optimization problem (Eq. 4) and the final decision function only require the computation of inner products (dot products) between the data points, $x_i^T x_j$. We never actually need to compute $w$ explicitly.
    \item \textbf{The Kernel Trick:} Because the algorithm only relies on inner products, we can seamlessly replace $x_i^T x_j$ with a Kernel function $K(x_i, x_j)$. This allows the algorithm to implicitly project the data into infinite-dimensional spaces to find linear separations for non-linear data, all while remaining computationally efficient.
\end{enumerate}

\subsection{The Role of Support Vectors}
\subsubsection{Why Only Data Points on the ``Margin'' Contribute to \texorpdfstring{$w$}{w}}
This is entirely due to the \textbf{Karush-Kuhn-Tucker (KKT) conditions}, specifically the condition of \textbf{complementary slackness}. 
In the optimal solution, the KKT complementary slackness condition states that for every data point $i$:
\begin{equation}
    \alpha_i \left( 1 - y_i(w^T x_i + b) \right) = 0
\end{equation}
This equation dictates that one of two things must be true for every data point:
\begin{itemize}
    \item \textbf{$y_i(w^T x_i + b) > 1$}: The point lies strictly outside the margin (it is correctly classified and safely away from the boundary). For the product in Equation 1 to be zero, the Lagrangian multiplier \textbf{$\alpha_i$ must equal $0$}.
    \item \textbf{$y_i(w^T x_i + b) = 1$}: The point lies exactly \textit{on} the margin. In this case, \textbf{$\alpha_i > 0$} is allowed.
\end{itemize}
Since the weight vector is defined as:
\begin{equation}
    w = \sum_{i=1}^n \alpha_i y_i x_i
\end{equation}
any point with $\alpha_i = 0$ vanishes entirely from the summation. Therefore, only the points where $\alpha_i > 0$ (the points sitting exactly on the margin, known as \textbf{Support Vectors}) play any role in defining the weight vector $w$.

\subsection{Simple SVM by Hand}
\subsubsection{For what range of parameter \texorpdfstring{$h > 0$}{h > 0}, are the training points still linearly separable?}
For two classes to be linearly separable, their bounding box\(conve\) of the data must not intersect. 
\begin{itemize}
    \item \textbf{Positive bounding box:} The line segment connecting $P_1(0,0)$ and $P_2(2,2)$. This sits on the line $y = x$ for $x \in [0, 2]$.
    \item \textbf{Negative bounding box:} The line segment connecting $N_1(h,1)$ and $N_2(0,3)$. The slope of this line is $\frac{3-1}{0-h} = -\frac{2}{h}$. The equation of this line is $y = -\frac{2}{h}x + 3$ for $x \in [0, h]$.
\end{itemize}
Let's find exactly where these two lines intersect by setting them equal:
\begin{align*}
    x &= -\frac{2}{h}x + 3 \\
    x \left(1 + \frac{2}{h}\right) &= 3 \implies x = \frac{3h}{h + 2}
\end{align*}
For the datasets to overlap (and thus \textit{not} be linearly separable), this intersection point must lie exactly on both line segments simultaneously.
\begin{enumerate}
    \item \textbf{Must be on the positive segment:} $0 \le x \le 2$
    \[ \frac{3h}{h + 2} \le 2 \implies 3h \le 2h + 4 \implies h \le 4 \]
    \item \textbf{Must be on the negative segment:} $0 \le x \le h$
    \[ \frac{3h}{h + 2} \le h \implies 3 \le h + 2 \implies h \ge 1 \]
\end{enumerate}
The bounding boxes intersect when $1 \le h \le 4$. Therefore, the points are \textbf{linearly separable} when:
\[ 0 < h < 1 \quad \textbf{OR} \quad h > 4 \]
\subsubsection{Does the orientation of the maximum margin decision boundary change as \texorpdfstring{$h$}{h} changes, when the points are separable?}
\subsubsection*{Conclusion:} It depends on the range. The orientation remains constant for $0 < h < 1$, but changes for $h > 4$. 
\subsubsection*{Explanation:}
\begin{itemize}
    \item \textbf{Case 1: $0 < h < 1$} \\
    In this range, the negative points are $(h,1)$ and $(0,3)$, both positioned to the ``left/above'' of the positive segment $y=x$. The support vectors will be $P_1(0,0)$, $P_2(2,2)$, and $N_1(h,1)$. Because the projection of $N_1$ onto the line segment connecting $P_1$ and $P_2$ falls strictly between them, the maximum margin boundary is tightly constrained by the parallel segment $P_1P_2$. As $h$ moves between $0$ and $1$, the boundary moves closer or further away, but its slope remains securely parallel to $y=x$ (slope = $1$). \textbf{The orientation does not change.}    
    \item \textbf{Case 2: $h > 4$} \\
    In this range, $N_1(h,1)$ is pulled far out to the right. The positive points are now ``below'' $N_2(0,3)$ and to the ``left'' of $N_1(h,1)$. The support vectors shift to $P_2(2,2)$, $N_2(0,3)$, and $N_1(h,1)$. As $h$ increases, the position of the support vector $N_1$ shifts horizontally. Because the boundary must thread the needle to maintain maximum distance between $P_2$ and the expanding negative hull $\{N_1, N_2\}$, the hyperplane must tilt. \textbf{The orientation does change.}
\end{itemize}

\section{Neural Networks and Backpropagation}
\textbf{Notation \& Definitions:}
Defining the variables for clarity as below:
\begin{itemize}
    \item \textbf{Input Data:} $x_i$ is the $i$-th input vector. $y_i$ is the target value.
    \item \textbf{Hidden Layer Weights:} $\theta_1$ and $\theta_2$ (replacing $\alpha$ and $\beta$).
    \item \textbf{Output Layer Weights:} $\phi = \begin{bmatrix} \phi_1 \\ \phi_2 \end{bmatrix}$ (replacing $w$).
    \item \textbf{Hidden Pre-activations:} $z_{i,1}^{(1)} = \theta_1^T x_i$ and $z_{i,2}^{(1)} = \theta_2^T x_i$.
    \item \textbf{Hidden Activations:} $a_i^{(1)} = \begin{bmatrix} a_{i,1}^{(1)} \\ a_{i,2}^{(1)} \end{bmatrix} = \begin{bmatrix} \sigma(z_{i,1}^{(1)}) \\ \sigma(z_{i,2}^{(1)}) \end{bmatrix}$ (replacing $z_i$).
    \item \textbf{Output Pre-activation:} $z_i^{(2)} = \phi^T a_i^{(1)}$ (replacing $u_i$).
    \item \textbf{Final Prediction:} $p_i = \sigma(z_i^{(2)})$ (replacing $\hat{y}_i$).
    \item \textbf{Cost Function:} The loss for a single sample is $L_i = (y_i - p_i)^2$. The total error is $\mathcal{E} = \sum_{i=1}^m L_i$.
\end{itemize}
\textbf{Standard Sigmoid Derivative Detail:} Recall that for $f(x) = \frac{1}{1 + e^{-x}}$, the derivative is calculated via the quotient/chain rule:
\[ \frac{d}{dx} \sigma(x) = \frac{d}{dx} (1 + e^{-x})^{-1} = -1(1 + e^{-x})^{-2} (-e^{-x}) = \frac{e^{-x}}{(1 + e^{-x})^2} \]
By adding and subtracting $1$ in the numerator, this famously simplifies to $\sigma(x)(1 - \sigma(x))$.

\subsection{Gradient with respect to the output weights \texorpdfstring{$\phi$}{phi}:}
We need to compute $\frac{\partial \mathcal{E}}{\partial \phi}$. By linearity of differentiation, this is the sum of the individual gradients $\frac{\partial L_i}{\partial \phi}$.
Applying the chain rule backward from the loss to the weights:
\begin{equation}
    \frac{\partial L_i}{\partial \phi} = \frac{\partial L_i}{\partial p_i} \cdot \frac{\partial p_i}{\partial z_i^{(2)}} \cdot \frac{\partial z_i^{(2)}}{\partial \phi}
\end{equation}
Let's break down these three components with explicit steps:\\
\textbf{Step 1a: Derivative of the loss w.r.t.\ the prediction ($p_i$)} \\
Apply the power rule and the chain rule (multiplying by the derivative of the inside, which is $-1$):
\[ \frac{\partial L_i}{\partial p_i} = \frac{\partial}{\partial p_i} (y_i - p_i)^2 = 2(y_i - p_i) \cdot \frac{\partial}{\partial p_i} (y_i - p_i) = 2(y_i - p_i)(-1) = -2(y_i - p_i) \]
\textbf{Step 1b: Derivative of the prediction with respect to the pre-activation ($z_i^{(2)}$)} \\
Using our sigmoid derivative identity:
\[ \frac{\partial p_i}{\partial z_i^{(2)}} = \sigma'(z_i^{(2)}) = \sigma(z_i^{(2)})(1 - \sigma(z_i^{(2)})) = p_i(1 - p_i) \]
\textit{Note: In backpropagation, we often combine 1a and 1b into a single ``Error Signal'' at the output layer, denoted as $\delta_i^{(2)} = -2(y_i - p_i)p_i(1 - p_i)$.}\\
\textbf{Step 1c: Derivative of the pre-activation w.r.t. the weights ($\phi$)} \\
Expanding the dot product, $z_i^{(2)} = \phi_1 a_{i,1}^{(1)} + \phi_2 a_{i,2}^{(1)}$. Differentiating with respect to the vector $\phi$ yields the activation vector from the previous layer:
\[ \frac{\partial z_i^{(2)}}{\partial \phi} = a_i^{(1)} \]
\textbf{Step 1d: Combine and sum} \\
Multiplying the terms, the gradient for a single example is:
\[ \frac{\partial L_i}{\partial \phi} = \underbrace{-2(y_i - p_i) p_i(1 - p_i)}_{\delta_i^{(2)}} \cdot a_i^{(1)} \]
Summing over the entire dataset of size $m$:
\begin{equation}
    \frac{\partial \mathcal{E}}{\partial \phi} = \sum_{i=1}^m \left[ -2(y_i - p_i) p_i(1 - p_i) a_i^{(1)} \right]
\end{equation}

\subsection{Deriving the gradient for \texorpdfstring{$\theta_1$}{theta1} (first hidden node):}
The weight vector $\theta_1$ exclusively controls the pre-activation $z_{i,1}^{(1)}$. The extended chain rule is:
\begin{equation}
    \frac{\partial L_i}{\partial \theta_1} = \left( \frac{\partial L_i}{\partial p_i} \cdot \frac{\partial p_i}{\partial z_i^{(2)}} \right) \cdot \frac{\partial z_i^{(2)}}{\partial a_{i,1}^{(1)}} \cdot \frac{\partial a_{i,1}^{(1)}}{\partial z_{i,1}^{(1)}} \cdot \frac{\partial z_{i,1}^{(1)}}{\partial \theta_1}
\end{equation}
We can reuse the error signal $\delta_i^{(2)}$ from the output layer. Now, let's derive the remaining three local gradients:
\begin{enumerate}
    \item \textbf{Backwards through the output weight:} How does the first hidden activation affect the output pre-activation?
    \[ \frac{\partial z_i^{(2)}}{\partial a_{i,1}^{(1)}} = \frac{\partial}{\partial a_{i,1}^{(1)}} (\phi_1 a_{i,1}^{(1)} + \phi_2 a_{i,2}^{(1)}) = \phi_1 \]
    \item \textbf{Backwards through the hidden sigmoid:}
    \[ \frac{\partial a_{i,1}^{(1)}}{\partial z_{i,1}^{(1)}} = \sigma'(z_{i,1}^{(1)}) = a_{i,1}^{(1)}(1 - a_{i,1}^{(1)}) \]
    \textit{Note: Combining $\delta_i^{(2)} \cdot \phi_1 \cdot a_{i,1}^{(1)}(1 - a_{i,1}^{(1)})$ gives us the hidden layer error signal, $\delta_{i,1}^{(1)}$.}
    \item \textbf{Backwards to the input data:}
    \[ \frac{\partial z_{i,1}^{(1)}}{\partial \theta_1} = \frac{\partial}{\partial \theta_1} (\theta_1^T x_i) = x_i \]
\end{enumerate}
\textbf{Total Gradient for $\theta_1$:}
\begin{equation}
    \frac{\partial \mathcal{E}}{\partial \theta_1} = \sum_{i=1}^m \left[ \underbrace{-2(y_i - p_i) p_i(1 - p_i)}_{\delta_i^{(2)}} \cdot \phi_1 \cdot a_{i,1}^{(1)}(1 - a_{i,1}^{(1)}) \cdot x_i \right]
\end{equation}

\subsection{Deriving the gradient for \texorpdfstring{$\theta_2$}{theta2} (second hidden node):}
The derivation for $\theta_2$ mirrors $\theta_1$ perfectly. The signal routes through the second hidden node $a_{i,2}^{(1)}$, meaning it scales by the second output weight $\phi_2$ instead of $\phi_1$.
The chain rule path is:
\begin{equation}
    \frac{\partial L_i}{\partial \theta_2} = \delta_i^{(2)} \cdot \frac{\partial z_i^{(2)}}{\partial a_{i,2}^{(1)}} \cdot \frac{\partial a_{i,2}^{(1)}}{\partial z_{i,2}^{(1)}} \cdot \frac{\partial z_{i,2}^{(1)}}{\partial \theta_2}
\end{equation}
Evaluating the local derivatives yields $\phi_2$, $a_{i,2}^{(1)}(1 - a_{i,2}^{(1)})$, and $x_i$, respectively.
\textbf{Total Gradient for $\theta_2$:}
\begin{equation}
    \frac{\partial \mathcal{E}}{\partial \theta_2} = \sum_{i=1}^m \left[ -2(y_i - p_i) p_i(1 - p_i) \cdot \phi_2 \cdot a_{i,2}^{(1)}(1 - a_{i,2}^{(1)}) \cdot x_i \right]
\end{equation}

\section{Feature Selection and Change-point Detection}

\subsection{Mutual Information based Feature Selection}
To determine which keyword is more informative, we calculate the Mutual Information (MI) for both ``prize'' and ``hello'' with respect to the email class ``spam''. 
The Mutual Information between two discrete random variables $X$ (keyword presence) and $Y$ (spam status) is given by the formula:
\begin{equation}
    I(X;Y) = \sum_{x \in \{0,1\}} \sum_{y \in \{0,1\}} p(x,y) \log_2 \left( \frac{p(x,y)}{p(x)p(y)} \right)
\end{equation}
First, let's establish the probabilities for each table by dividing the counts by the total number of emails in each respective dataset.
\subsubsection{Mutual Information for ``prize'' \texorpdfstring{($X$)}{(X)} and ``spam'' \texorpdfstring{($Y$)}{(Y)}}
Total emails ($N$) = $145 + 35 + 1540 + 16740 = \textbf{18,460}$
\subsubsection{Marginal Probabilities:}
\begin{align*}
    P(Y=1) &= (145 + 35) / 18460 = 180 / 18460 \approx 0.00975 \\
    P(Y=0) &= (1540 + 16740) / 18460 = 18280 / 18460 \approx 0.99025 \\
    P(X=1) &= (145 + 1540) / 18460 = 1685 / 18460 \approx 0.09128 \\
    P(X=0) &= (35 + 16740) / 18460 = 16775 / 18460 \approx 0.90872
\end{align*}
\subsubsection{Joint Probabilities \& MI Terms:}
\begin{itemize}
    \item \textbf{$X=1, Y=1$:} $p(1,1) = 145 / 18460 \approx 0.00785$ \\
    $\text{Term} = 0.00785 \times \log_2 \left( \frac{0.00785}{0.09128 \times 0.00975} \right) \approx 0.02468$
    \item \textbf{$X=0, Y=1$:} $p(0,1) = 35 / 18460 \approx 0.00190$ \\
    $\text{Term} = 0.00190 \times \log_2 \left( \frac{0.00190}{0.90872 \times 0.00975} \right) \approx -0.00422$
    \item \textbf{$X=1, Y=0$:} $p(1,0) = 1540 / 18460 \approx 0.08342$ \\
    $\text{Term} = 0.08342 \times \log_2 \left( \frac{0.08342}{0.09128 \times 0.99025} \right) \approx -0.00965$
    \item \textbf{$X=0, Y=0$:} $p(0,0) = 16740 / 18460 \approx 0.90683$ \\
    $\text{Term} = 0.90683 \times \log_2 \left( \frac{0.90683}{0.90872 \times 0.99025} \right) \approx 0.00972$
\end{itemize}
\textbf{Total MI for ``prize'':} $0.02468 - 0.00422 - 0.00965 + 0.00972 \approx \textbf{0.0205 bits}$
\subsubsection{Mutual Information for ``hello'' \texorpdfstring{($X$)}{(X)} and ``spam'' \texorpdfstring{($Y$)}{(Y)}}
Total emails ($N$) = $190 + 55 + 12200 + 6740 = \textbf{19,185}$
\subsubsection{Marginal Probabilities:}
\begin{align*}
    P(Y=1) &= 245 / 19185 \approx 0.01277 \\
    P(Y=0) &= 18940 / 19185 \approx 0.98723 \\
    P(X=1) &= 12390 / 19185 \approx 0.64582 \\
    P(X=0) &= 6795 / 19185 \approx 0.35418
\end{align*}
\subsubsection{Joint Probabilities \& MI Terms:}
\begin{itemize}
    \item \textbf{$X=1, Y=1$:} $p(1,1) = 190 / 19185 \approx 0.00990$ \\
    $\text{Term} = 0.00990 \times \log_2 \left( \frac{0.00990}{0.64582 \times 0.01277} \right) \approx 0.00262$
    \item \textbf{$X=0, Y=1$:} $p(0,1) = 55 / 19185 \approx 0.00287$ \\
    $\text{Term} = 0.00287 \times \log_2 \left( \frac{0.00287}{0.35418 \times 0.01277} \right) \approx -0.00189$
    \item \textbf{$X=1, Y=0$:} $p(1,0) = 12200 / 19185 \approx 0.63591$ \\
    $\text{Term} = 0.63591 \times \log_2 \left( \frac{0.63591}{0.64582 \times 0.98723} \right) \approx -0.00254$
    \item \textbf{$X=0, Y=0$:} $p(0,0) = 6740 / 19185 \approx 0.35132$ \\
    $\text{Term} = 0.35132 \times \log_2 \left( \frac{0.35132}{0.35418 \times 0.98723} \right) \approx 0.00242$
\end{itemize}
\textbf{Total MI for ``hello'':} $0.00262 - 0.00189 - 0.00254 + 0.00242 \approx \textbf{0.0006 bits}$
\subsubsection{Python Code Used for Verification:}
\begin{lstlisting}
import numpy as np
def calc_mi(c_11, c_10, c_01, c_00):
    total = c_11 + c_10 + c_01 + c_00
    p11, p10 = c_11/total, c_10/total
    p01, p00 = c_01/total, c_00/total
    
    px1, px0 = p11 + p01, p10 + p00
    py1, py0 = p11 + p10, p01 + p00
    
    mi = 0
    if p11 > 0: mi += p11 * np.log2(p11 / (px1 * py1))
    if p10 > 0: mi += p10 * np.log2(p10 / (px0 * py1))
    if p01 > 0: mi += p01 * np.log2(p01 / (px1 * py0))
    if p00 > 0: mi += p00 * np.log2(p00 / (px0 * py0))
    return mi
print(f"MI for 'prize': {calc_mi(145, 35, 1540, 16740)}")
print(f"MI for 'hello': {calc_mi(190, 55, 12200, 6740)}")
\end{lstlisting}
\subsubsection{Conclusion:}
The keyword \textbf{``prize''} is significantly more informative ($0.0205$ bits) compared to the keyword ``hello'' ($0.0006$ bits) for deciding whether an email is spam.
\subsection{Change-point Detection: CUSUM Statistic}
Given the normal distributions $f_0 = \mathcal{N}(0, 1)$ and $f_1 = \mathcal{N}(0, 1.25)$ where the variance changes from $1$ to $1.25$, the Probability Density Functions (PDF) are:
\begin{align*}
    f_0(x) &= \frac{1}{\sqrt{2\pi}} \exp\left(-\frac{x^2}{2}\right) \\
    f_1(x) &= \frac{1}{\sqrt{2\pi(1.25)}} \exp\left(-\frac{x^2}{2(1.25)}\right)
\end{align*}
The CUSUM method tracks the cumulative sum of the log-likelihood ratio. For a single observation $x_i$, the log-likelihood ratio $W_i$ is derived as:
\begin{equation*}
    W_i = \ln \left( \frac{f_1(x_i)}{f_0(x_i)} \right) = \ln \left( \frac{\frac{1}{\sqrt{2\pi(1.25)}} \exp\left(-\frac{x_i^2}{2.5}\right)}{\frac{1}{\sqrt{2\pi}} \exp\left(-\frac{x_i^2}{2}\right)} \right)
\end{equation*}
Simplifying the natural log:
\begin{align*}
    W_i &= \ln \left( \frac{1}{\sqrt{1.25}} \right) - \frac{x_i^2}{2.5} + \frac{x_i^2}{2} \\
    W_i &= -0.5 \ln(1.25) + x_i^2 \left(\frac{1}{2} - \frac{1}{2.5}\right) \\
    W_i &= 0.1 x_i^2 - 0.5 \ln(1.25) \approx 0.1 x_i^2 - 0.11157
\end{align*}
The \textbf{CUSUM Statistic ($S_t$)} is updated recursively as:
\begin{align*}
    S_0 &= 0 \\
    S_t &= \max(0, S_{t-1} + 0.1 x_t^2 - 0.5 \ln(1.25))
\end{align*}
\subsubsection{Plotting and Change Estimation}
Below is the Python code to generate the $150$ samples and plot the CUSUM statistic, utilizing the required random seed ($6740$).
\begin{lstlisting}
import numpy as np
import matplotlib.pyplot as plt

# 1. Setup and Seed
np.random.seed(6740)

# 2. Generate Data
# First 100 samples from N(0, 1)
x_f0 = np.random.normal(0, np.sqrt(1.0), 100)
# Next 50 samples from N(0, 1.25)
x_f1 = np.random.normal(0, np.sqrt(1.25), 50)
x = np.concatenate((x_f0, x_f1))

# 3. Calculate CUSUM Statistic
S = np.zeros(len(x) + 1)
threshold_constant = 0.5 * np.log(1.25)

for t in range(1, len(x) + 1):
    W_t = 0.1 * (x[t-1]**2) - threshold_constant
    S[t] = max(0, S[t-1] + W_t)

# 4. Plotting
plt.figure(figsize=(10, 5))
plt.plot(range(len(S)), S, label='CUSUM Statistic (S_t)')
plt.axvline(x=100, color='r', linestyle='--', label='True Change Point (t=100)')
plt.title('CUSUM Statistic for Change in Variance')
plt.xlabel('Time (t)')
plt.ylabel('S_t')
plt.legend()
plt.grid(True)
plt.savefig('cusum_change_point_detection.png', bbox_inches='tight')
\end{lstlisting}
\subsubsection{Reasonable Estimation of the Change-Point:}
\begin{figure}[H]
    \centering
    \includegraphics[width=0.85\textwidth]{ISYE6740_Summer_2026_HW4/cusum_change_point_detection.png}
    \caption{CUSUM Statistic tracking a change in variance. The sustained upward trajectory begins just after $t=100$.}\label{fig:cusum}
\end{figure}
Running the script above, we observe the CUSUM statistic $S_t$ staying relatively flat and repeatedly resetting to zero during the first $100$ samples (as the expected value of $W_t$ under $f_0$ is negative). 
Shortly after sample index $100$, the statistic begins a persistent, sustained upward trajectory because the expected value of $W_t$ under $f_1$ becomes positive due to the higher variance producing larger $x_i^2$ values. Therefore, a reasonable estimation for the detected change point from the plot is precisely where the final upward slope firmly begins (which practically materializes slightly after $t = 100$).

\end{document}